% !TEX program = lualatex
% Transparencias de Fundamentos de las Matemáticas II
% Clase 11: Aritmética modular III

\documentclass[aspectratio=169,11pt]{beamer}

\usepackage{fontspec}
\usepackage[spanish,es-nodecimaldot]{babel}
\usepackage{amsmath,amssymb,mathtools}
\usepackage{booktabs}
\usepackage{csquotes}
\usepackage{xcolor}

\usetheme{Madrid}
\usecolortheme{default}
\setbeamertemplate{navigation symbols}{}
\setbeamertemplate{footline}[frame number]

\definecolor{FdMBlue}{RGB}{20,55,100}
\definecolor{FdMRed}{RGB}{130,30,30}
\definecolor{FdMGreen}{RGB}{25,100,60}

\setbeamercolor{structure}{fg=FdMBlue}
\setbeamercolor{title}{fg=white,bg=FdMBlue}
\setbeamercolor{frametitle}{fg=white,bg=FdMBlue}
\setbeamercolor{block title}{fg=white,bg=FdMBlue}
\setbeamercolor{block body}{bg=blue!3}
\setbeamercolor{alerted text}{fg=FdMRed}

\setmainfont{Latin Modern Roman}
\setsansfont{Latin Modern Sans}
\setmonofont{Latin Modern Mono}

\newcommand{\N}{\mathbb{N}}
\newcommand{\Z}{\mathbb{Z}}
\newcommand{\Q}{\mathbb{Q}}
\newcommand{\R}{\mathbb{R}}
\newcommand{\C}{\mathbb{C}}
\newcommand{\Pow}{\mathcal{P}}
\newcommand{\id}{\operatorname{id}}
\newcommand{\mcd}{\operatorname{mcd}}
\newcommand{\mcm}{\operatorname{mcm}}
\newcommand{\im}{\operatorname{Im}}

\title[FdM II -- Clase 11]{Aritmética modular III}
\subtitle{Congruencias lineales y Teorema Chino del Resto}
\author{Antonio Falcó Montesinos}
\institute{Universidad CEU Cardenal Herrera}
\date{Fundamentos de las Matemáticas II}

\begin{document}

\begin{frame}
  \titlepage
\end{frame}

\begin{frame}{Objetivos de la clase}
\begin{itemize}
  \item Resolver congruencias lineales.
  \item Aplicar el criterio \(\mcd(a,n)mid b\).
  \item Enunciar y usar el Teorema Chino del Resto.
\end{itemize}
\end{frame}

\begin{frame}{Mapa de la clase}
\tableofcontents
\end{frame}

\section{Congruencias lineales}

\begin{frame}{Congruencias lineales}
\begin{block}{Forma}
\begin{itemize}
  \item Una congruencia lineal tiene la forma \(ax\equiv b\pmod n\).
  \item La incógnita es un entero \(x\).
  \item Se buscan clases de soluciones módulo \(n\).
\end{itemize}
\end{block}
\end{frame}

\begin{frame}{Congruencias lineales}
\begin{block}{Criterio}
\begin{itemize}
  \item \(ax\equiv b\pmod n\) tiene solución si y solo si \(\mcd(a,n)mid b\).
  \item Si \(\mcd(a,n)=1\), se puede multiplicar por el inverso de \([a]_n\).
  \item Si el \(\mcd\) es mayor que \(1\), hay que reducir con cuidado.
\end{itemize}
\end{block}
\end{frame}

\begin{frame}{Congruencias lineales}
\begin{block}{Ejemplo}
\begin{itemize}
  \item Resolver \(3x\equiv 4\pmod 7\).
  \item El inverso de \(3\) módulo \(7\) es \(5\).
  \item Entonces \(x\equiv 5\cdot 4\equiv 6\pmod 7\).
\end{itemize}
\end{block}
\end{frame}

\section{Teorema Chino}

\begin{frame}{Teorema Chino}
\begin{block}{Enunciado}
\begin{itemize}
  \item Si los módulos \(n_1,\ldots,n_r\) son coprimos dos a dos, el sistema tiene solución única módulo \(N=n_1\cdots n_r\).
  \item El teorema permite reconstruir una clase a partir de varios restos.
  \item La coprimalidad dos a dos es la hipótesis clave.
\end{itemize}
\end{block}
\end{frame}

\begin{frame}{Teorema Chino}
\begin{block}{Ejemplo}
\begin{itemize}
  \item Resolver \(x\equiv 2\pmod 3\), \(x\equiv 3\pmod 5\).
  \item Probando clases módulo \(15\), se obtiene \(x\equiv 8\pmod {15}\).
  \item Todas las soluciones son \(x=8+15k\).
\end{itemize}
\end{block}
\end{frame}

\section{Profundización}

\begin{frame}{Profundización}
\begin{block}{Dividir congruencias}
\begin{itemize}
  \item No siempre se puede dividir por un entero.
  \item Dividir por \(d\) exige modificar también el módulo si \(d\) divide a todos los términos.
  \item La forma segura es usar el criterio del \(\mcd\).
\end{itemize}
\end{block}
\end{frame}

\begin{frame}{Profundización}
\begin{block}{Construcción TCR}
\begin{itemize}
  \item Se trabaja módulo \(N=n_1\cdots n_r\).
  \item Para cada módulo se construye un término que valga \(1\) en ese módulo y \(0\) en los demás.
  \item La suma ponderada produce la solución.
\end{itemize}
\end{block}
\end{frame}

\begin{frame}{Profundización}
\begin{block}{Errores frecuentes}
\begin{itemize}
  \item Aplicar TCR con módulos no coprimos.
  \item Dar una solución entera y olvidar la clase módulo \(N\).
  \item Cancelar sin comprobar invertibilidad.
\end{itemize}
\end{block}
\end{frame}


\section{Detalle y práctica}

\begin{frame}{Detalle y práctica}
\begin{block}{Inversos modulares}
\begin{itemize}
  \item Un elemento \(a\) tiene inverso módulo \(n\) si existe \(x\) tal que \(ax\equiv 1\pmod n\).
  \item Esto ocurre exactamente cuando \(\mcd(a,n)=1\).
  \item El algoritmo extendido de Euclides construye el inverso.
\end{itemize}
\end{block}
\end{frame}

\begin{frame}{Detalle y práctica}
\begin{block}{Actividad breve}
\begin{itemize}
  \item Encontrar el inverso de \(7\) módulo \(26\).
  \item Resolver \(7x\equiv 5\pmod {26}\).
  \item Explicar por qué \(6x\equiv 1\pmod {26}\) no tiene solución.
\end{itemize}
\end{block}
\end{frame}

\begin{frame}{Cierre}
\begin{itemize}
  \item Las congruencias lineales se deciden con el \(\mcd\).
  \item El TCR combina restos compatibles.
  \item La unicidad siempre es módulo el producto de los módulos.
\end{itemize}
\end{frame}

\begin{frame}{Trabajo recomendado}
\begin{itemize}
  \item Releer las secciones correspondientes del manual.
  \item Identificar definiciones, símbolos nuevos y enunciados principales.
  \item Preparar dudas y ejemplos para el seminario asociado.
\end{itemize}
\end{frame}

\end{document}
